% !PRINTED 106
\weberpage{106}
%<<CONT
\[ \mathfrak{A}_3)\quad \eta = i\,\frac{\sqrt{3}-1}{\sqrt{2}} \]
\[ \mathfrak{A}_4)\quad \eta = \frac{1+\sqrt{3}}{\sqrt{2}} \]
\[ \mathfrak{A}_5)\quad \eta = \frac{1+i\sqrt{3}}{\sqrt{2}}. \]

Finally, the case of a single invariant $i_1 = 4$ shall now
be considered. The conditions for $m$ then consist simply in
this, that $m$ can be divisible by 4, 8, 16, but by no higher
power of 2, that odd primes can divide $m$ only once, and that
at least one of the index moduli must be divisible by 4. The
first values are thus
\[ m = 5,\ 13,\ 15,\ 16,\ 17,\ 20,\ 29,\ 33,\ 37,\ 39,\ 40\ \ldots \]

We shall assume that of the index moduli,
$c_1, c_2, \ldots c_\varrho$ are divisible by 4, and $c_{\varrho+1}, c_{\varrho+2}, \ldots c_\mu$ by 2, but
not by 4. Then $\varrho$ must be at least equal to 1,
and consequently
\[ \begin{aligned}
&\delta_{1,1} = \delta_{1,2} = \cdots = \delta_{1,\varrho} = 4, \quad \delta_{1,\varrho+1} = \cdots = \delta_{1,\mu} = 2\\
&i_{1,1} = i_{1,2} = \cdots = i_{1,\varrho} = 1, \quad i_{1,\varrho+1} = \cdots = i_{1,\mu} = 2\\
&c_{1,1} = \frac{c_1}{4}, \cdots\cdots, c_{1,\varrho} = \frac{c_\varrho}{4}, \quad c_{1,\varrho+1} = \frac{c_{\varrho+1}}{2}, \cdots, c_{1,\mu} = \frac{c_\mu}{2}.
\end{aligned} \]

Now the $e_{1,h}$ are to be determined so that of the numbers
\[ e_{1,1}, e_{1,2}, \ldots e_{1,\varrho} \]
at least one is odd, and that (on account of \S.\,21, 6.) of the
numbers
\[ e_{1,1}, e_{1,2}, \ldots e_{1,\varrho} \]
none is divisible by 4, and of the
\[ e_{1,\varrho+1}, e_{1,\varrho+2}, \ldots e_{1,\mu} \]
none is divisible by 2. An exception is again formed here,
in the case where $m$ is divisible by 8 or by 16, by the number $e_{1,h}$
which corresponds to none of the prime factors of $m$, and which can also
be even.

Then for the indices $\alpha$ of the numbers $a$ we obtain the congruence:
\[ e_{1,1}\alpha_1 + \cdots + e_{1,\varrho}\alpha_\varrho + 2\,(e_{1,\varrho+1}\alpha_{\varrho+1} + \cdots + e_{1,\mu}\alpha_\mu) \equiv 0\ (\mathrm{mod}\ 4). \]
